\input{../common}
\input{../bibliography}

\begin{document}
  %<*content>
  \sheet{algebra}{invariants-de-similitude}{Invariants de similitude}

  \summary{Nous montrons l'existence et l'unicité des invariants de similitude d'un endomorphisme d'un espace de dimension finie en utilisant la dualité.}

  \reference[GOU21]{398}

  Soit $E$ un espace vectoriel de dimension finie $n \geq 1$ sur un corps commutatif $\mathbb{K}$. Soit $f \in \mathcal{L}(E)$.

  \begin{notation}
    Soit $x \in E$. On note $P_x$ le polynôme unitaire engendrant l'idéal $\{ P \in \mathbb{K}[X] \mid P(f)(x) = 0 \}$ (un tel polynôme existe car $\mathbb{K}[X]$ est principal et cet idéal est non réduit à $\{ 0 \}$) et $E_x = \{ P(f)(x) \mid P \in \mathbb{K}[X] \}$.
  \end{notation}

  \begin{lemma}
    \label{invariants-de-similitude-1}
    \begin{enumerate}[label={\bfseries(\roman*)}]
      \item Si $k = \deg(\pi_f)$, alors $\mathbb{K}[f]$ est un sous-espace vectoriel de $\mathcal{L}(E)$ de dimension $k$, dont une base est $(f^i)_{i \in \llbracket 0, k-1 \rrbracket}$.
      \item Soit $x \in E$. Si $l = \deg(P_x)$, alors $E_x$ est un sous-espace vectoriel de $E$ de dimension $l$, dont une base est $(f^i(x))_{i \in \llbracket 0, l-1 \rrbracket}$.
    \end{enumerate}
  \end{lemma}

  \reference{61}

  \begin{proof}
    \begin{enumerate}[label={\bfseries(\roman*)}]
      \item Montrons que la famille $(f^i)_{i \in \llbracket 0, k-1 \rrbracket}$ est à la fois libre et génératrice.
      \begin{itemize}
        \item Soit $P(f) \in \mathbb{K}[f]$. On fait la division euclidienne de $P$ par $\pi_f$ dans $\mathbb{K}[X]$ pour écrire $P = \pi_f Q + R$ avec $Q, R \in \mathbb{K}[X]$ et $\deg(R) < k = \deg(\pi_f)$. En évaluant en $f$, cela donne $P(f) = R(f) \in \operatorname{Vect}(\id_E, \dots, f^{k-1})$. Donc la famille est génératrice.
        \item Si $\sum_{i=0}^{k-1} \lambda_i f^i = 0$, alors le polynôme $P = \sum_{i=0}^{k-1} \lambda_i X^i$ vérifie $P(f) = 0$. Donc $\pi_f \mid P$, et comme $\deg(P) < \deg(\pi_f)$, on a $P = 0$. Donc $\lambda_0 = \dots = \lambda_{k-1} = 0$. Donc la famille est libre.
      \end{itemize}
      \item La deuxième assertion se montre sensiblement de la même manière.
    \end{enumerate}
  \end{proof}

  \reference{290}

  \begin{lemma}
    \label{invariants-de-similitude-2}
    Il existe $x \in E$ tel que $P_x = \pi_f$.
  \end{lemma}

  La démonstration est un peu trop longue pour être incluse ici : c'est un résultat qui demande du temps pour le démontrer (et pourrait constituer un vrai développement à part entière). Nous vous renvoyons vers \cite{[GOU21]} p. 178 pour la démonstration.

  \begin{theorem}[Théorème de Frobenius]
    Il existe des sous-espaces vectoriels $F_1, \dots, F_r$ de $E$ tous stables par $f$ tels que :
    \begin{enumerate}[label={\bfseries(\roman*)}]
      \item \label{invariants-de-similitude-3} $E = \bigoplus_{i = 1}^r F_i$.
      \item \label{invariants-de-similitude-4} $\forall i \in \llbracket 1, r \rrbracket$, la restriction $f_i = f_{|F_i}$ est un endomorphisme cyclique de $F_i$.
      \item \label{invariants-de-similitude-5} Si $P_i = \pi_{f_i}$ est le polynôme minimal de $f_i$, on a $P_{i+1} \mid P_i$ $\forall i \in \llbracket 1, r-1 \rrbracket$.
    \end{enumerate}
    La suite $(P_i)_{i \in \llbracket 1, r \rrbracket}$ ne dépend que de $f$ et non du choix de la décomposition (elle est donc unique). On l'appelle \textbf{suite des invariants de $f$}.
  \end{theorem}

  \begin{proof}
    \begin{itemize}
      \item \uline{Existence :} Soit $k = \deg(\pi_f)$. Par le \cref{invariants-de-similitude-2}, il existe $x \in E$ tel que $P_x = \pi_f$. Par le \cref{invariants-de-similitude-1}, le sous-espace $F = E_x$ est de dimension $k$ et est stable par $f$ et comme $\deg(P_x) = k$, la famille de vecteurs
      \[ (\underbrace{x}_{= e_1}, \dots, \underbrace{f^{k-1}(x)}_{= e_k}) \]
      forme une base de $F$. Complétons cette base en une base $(e_1, \dots, e_n)$ de $E$. En désignant par $(e_1^*, \dots, e_n^*)$ la base duale associée et en notant $\Gamma = \{ e_k^* \circ f^i \mid i \in \mathbb{N} \}$, on pose
      \begin{align*}
        G &= \Gamma^\circ \\
        &= \{ x \in E \mid \forall i \in \mathbb{N}, \, (e_k^* \circ f^i)(x) = 0 \}
      \end{align*}
      Ainsi, $G$ est l'ensemble des $x \in E$ tels que la $k$-ième coordonnée de $f^i(x)$ (dans la base $(e_1, \dots, e_n)$) est nulle $\forall i \in \mathbb{N}$ ; $G$ est donc un sous-espace de $E$ stable par $f$. Montrons que $F \oplus G = E$.
      \newpar
      Montrons que $F \, \cap \, G = \{ 0 \}$. Soit $y \in F \, \cap \, G$. Si $y \neq 0$, on peut écrire $y = \lambda_1 e_1 + \dots + \lambda_p e_p$ avec $\lambda_p \neq 0$ et $p \leq k$. En composant par $e_k^* \circ f^{k-p}$, on obtient
      \begin{align*}
        0 &\underset{y \in G}{=} e_k^* \circ f^{k-p}(y) \\
        &= e_k^* (\lambda_1 f^{k-p}(e_1) + \dots + \lambda_p f^{k-p}(e_p)) \\
        &= e_k^* (\lambda_1 f^{k-p}(x) + \dots + \lambda_p f^{k-p}(x)) \\
        &= \lambda_p
      \end{align*}
      Ce qui est absurde.
      \newpar
      Montrons que $\dim(F) + \dim(G) = n$. Cela revient à montrer que $\dim(G) = n - k$. On sait que $G = \Gamma^\circ = (\operatorname{Vect}(\Gamma))^\circ$ et $\dim(\operatorname{Vect}(\Gamma)) + \dim(\operatorname{Vect}(\Gamma)^\circ) = n$. Montrons donc que $\dim(\operatorname{Vect}(\Gamma)) = k$. Posons
      \[
      \varphi :
      \begin{array}{ccc}
        \mathbb{K}[f] &\rightarrow& \operatorname{Vect}(\Gamma) \\
        g &\mapsto& e_k^* \circ g
      \end{array}
      \]
      Par définition de $\Gamma$, $\varphi$ est surjective. Soit $g \in \ker(\varphi)$. On a alors $e_k^* \circ g = 0$, et comme $g \in \mathbb{K}[f]$,
      \[ g = \lambda_1 \id + \dots \lambda_p f^{p-1} \text{ avec } \lambda_p \neq 0 \text{ et } p \leq k \]
      On a donc $0 = e_k^* \circ g(f^{k-p} (x)) = \lambda_p \neq 0$. Ainsi, $g = 0$ et $\phi$ est un isomorphisme. Donc $\dim(\operatorname{Vect}(\Gamma)) = \dim(\mathbb{K}[f]) = k$ par le \cref{invariants-de-similitude-1}, ce que l'on voulait.
      \newpar
      Soit $P_1$ le polynôme minimal de $f_{|F}$ (qui est le polynôme minimal de $f$ car $P_1 = \pi_{f_{|F}} = \underset{\pi_f = P_x}{=} \pi_f$). Soit $P_2$ le polynôme minimal de $f_{|G}$. Comme $G$ est stable par $f$, on a $P_1(f_{|G}) = \pi_f(f_{|G}) = 0$, donc $P_2 \mid P_1$. Il suffit alors de réitérer en remplaçant $f$ par $f_{|G}$ et $E$ par $G$ pour obtenir la décomposition voulu.
      \item \uline{Unicité :} Soient $F_1, \dots, F_r$ et $G_1, \dots G_s$ des sous-espaces vectoriels stables par $f$ qui vérifient le \cref{invariants-de-similitude-3}, le \cref{invariants-de-similitude-4} et le \cref{invariants-de-similitude-5}. On note pour tout $i$, $P_i = \pi_{f_{|F_i}}$ et $Q_i = \pi_{f_{|G_i}}$. On suppose par l'absurde $(P_1, \dots, P_r) \neq (Q_1, \dots, Q_s)$. Soit $j = \min\{ i \mid P_i \neq Q_i \}$.
      Comme $E = \bigoplus_{i = 1}^r F_i$ (où $\forall i \in \llbracket 1, r \rrbracket$, $F_i$ est stable par $f$ et $\forall k \geq j \geq 1$, $P_j(f)(F_k) = 0$) :
      \[ P_j(f)(F_1) \oplus \dots \oplus P_j(f)(F_{j-1}) = P_j(f)(E) \tag{$*$} \]
      De même,
      \[ P_j(f)(G_1) \oplus \dots \oplus P_j(f)(G_{j-1}) \oplus P_j(f)(G_j) \oplus \dots \oplus P_j(f)(G_s) = P_j(f)(E) \tag{$**$} \]
      Notons que l'on a $\forall i \in \llbracket 1, j-1 \rrbracket$, $\dim(P_j(f)(F_i)) = \dim(P_j(f)(G_i))$. En effet, on peut trouver une base $\mathcal{B}_i$ de $F_i$ et une base $\mathcal{B}'_i$ de $G_i$ telles que $\operatorname{Mat}(f_{|F_i}, \mathcal{B}_i) = \operatorname{Mat}(f_{|G_i}, \mathcal{B}'_i)$ par cyclicité de $f_{|F_i}$ et $f_{|G_i}$. En prenant les dimensions dans $(*)$ et $(**)$, on en déduit :
      \[ 0 = \dim(P_j(f)(G_j)) = \dots = \dim(P_j(f)(G_s)) \implies Q_j \mid P_j \]
      Par symétrie, on a de même $P_j \mid Q_j$. D'où $P_j = Q_j$ : absurde.
    \end{itemize}
  \end{proof}
  %</content>
\end{document}
